\DocumentMetadata{
  pdfversion=2.0,
  pdfstandard=ua-2,
  lang=en-US,
  tagging=on,
  tagging-setup={math/setup=mathml-SE,tabsorder=structure}
}
\documentclass[11pt,letterpaper]{article}
\usepackage[margin=0.68in,includefoot]{geometry}
\usepackage{newcomputermodern}
\usepackage{amsmath,amscd,mathtools}
\usepackage{tikz,adjustbox}
\providecommand{\square}{\mdlgwhtsquare}
\usepackage[hidelinks]{hyperref}
\hypersetup{
  pdftitle={Analysis PhD Exam, October 1994},
  pdfauthor={University of Florida Department of Mathematics},
  pdfsubject={Graduate mathematics examination},
  pdfdisplaydoctitle=true,
  bookmarksopen=true
}
\setcounter{secnumdepth}{0}
\setlength{\parindent}{0pt}
\setlength{\parskip}{0.28em}
\setlength{\emergencystretch}{3em}
\setlength{\leftmargini}{2.15em}
\setlength{\leftmarginii}{2.65em}
\setlength{\leftmarginiii}{3em}
\pagestyle{empty}
\raggedbottom
\allowdisplaybreaks
\definecolor{ProblemConcern}{HTML}{7A1F1F}
\newenvironment{problemconcern}{%
  \par\tagstructbegin{tag=Aside}\begingroup\color{ProblemConcern}
}{%
  \par\endgroup\tagstructend\nopagebreak[4]\vspace{0.12em}
}
\NewTaggingSocketPlug{block/list/label}{examproblem}{%
  \tagstructbegin{tag=Lbl}\tagstructend%
  \tagstructbegin{tag=\LBody}%
  \tagstructbegin{tag=\ProblemHeadingTag,title-o={\ProblemHeadingText}}%
  \tagmcbegin{tag=\ProblemHeadingTag,actualtext={\ProblemHeadingText}}%
  #2%
  \tagmcend%
  \tagstructend%
}
\newenvironment{examproblems}[2]{%
  \def\ProblemHeadingTag{#2}%
  \AssignTaggingSocketPlug{block/list/label}{examproblem}%
  \begin{enumerate}%
  \setcounter{enumi}{#1}%
  \renewcommand{\labelenumi}{\textbf{\theenumi.}}%
  \setlength{\topsep}{0.35em}%
  \setlength{\partopsep}{0pt}%
  \setlength{\itemsep}{0.48em}%
  \setlength{\parsep}{0.18em}%
}{\end{enumerate}}
\newenvironment{examsubparts}{%
  \AssignTaggingSocketPlug{block/list/label}{default}%
  \begin{enumerate}%
  \setlength{\topsep}{0.18em}%
  \setlength{\partopsep}{0pt}%
  \setlength{\itemsep}{0.16em}%
  \setlength{\parsep}{0.1em}%
}{\end{enumerate}}
\newenvironment{examinstructions}{%
  \par\begingroup\itshape
}{%
  \par\endgroup\vspace{0.18em}
}
\makeatletter
\renewcommand\subsection{\@startsection{subsection}{2}{0pt}%
  {-1.25ex plus -.25ex minus -.1ex}%
  {0.55ex plus .1ex}%
  {\normalfont\large\bfseries}}
\renewcommand{\@maketitle}{%
  \newpage\null\vskip -1em%
  {\footnotesize\bfseries
    \href{https://www.ufl.edu/}{University of Florida}, \href{https://math.ufl.edu/}{Department of Mathematics}\par}%
  \vskip -0.5em%
  \tagstructbegin{tag=H1,title={Analysis PhD Exam, October 1994}}%
  \tagpdfparaOff%
  \tagmcbegin{tag=H1}%
    {\LARGE\bfseries\href{https://gma.math.ufl.edu/exams/analysis-phd/}{Analysis PhD Exam}, October 1994\par}%
  \tagmcend%
  \tagpdfparaOn%
  \tagstructend%
  \vskip 0.25em%
  \tagmcbegin{artifact}%
    \hrule height 1pt%
  \tagmcend%
  \vskip 1em%
}
\makeatother
\title{Analysis PhD Exam, October 1994}
\author{}
\date{}
\begin{document}
\maketitle
\thispagestyle{empty}
\begin{examinstructions}
\((X,\Sigma,\mu)\) is a finite measure space.
\end{examinstructions}
\subsection{I.}
\begin{examinstructions}
State the following theorems:
\end{examinstructions}
\begin{examproblems}{0}{H3}
\def\ProblemHeadingText{Problem 1.}
\item
Lebesgue dominated convergence theorem
\def\ProblemHeadingText{Problem 2.}
\item
Fatou Lemma
\def\ProblemHeadingText{Problem 3.}
\item
Vitali theorem
\def\ProblemHeadingText{Problem 4.}
\item
Monotone convergence theorem
\def\ProblemHeadingText{Problem 5.}
\item
The Radon–Nikodym theorem
\def\ProblemHeadingText{Problem 6.}
\item
The Egorov theorem
\def\ProblemHeadingText{Problem 7.}
\item
The Fubini theorem
\def\ProblemHeadingText{Problem 8.}
\item
The Lebesgue decomposition theorem
\def\ProblemHeadingText{Problem 9.}
\item
The Uniform boundedness theorem
\def\ProblemHeadingText{Problem 10.}
\item
The Hahn–Banach theorem
\end{examproblems}
\subsection{II.}
\begin{examproblems}{0}{H3}
\def\ProblemHeadingText{Problem 1.}
\item
Prove the \((L^1)^*=L^\infty\)
\end{examproblems}
\subsection{III. Problems}
\begin{examproblems}{0}{H3}
\def\ProblemHeadingText{Problem 1.}
\item
Find \(\displaystyle \lim_{n\to\infty}\int_0^1 n x^n\,dx\). (Hint: use the monotone convergence theorem)
\def\ProblemHeadingText{Problem 2.}
\item
Let~\(\mu\) be the Lebesgue measure on~\(\mathbb{R}\) and~\(f\) a Lebesgue-integrable function on~\(\mathbb{R}\), such that \(\int_0^x f\,d\mu=0\) for every \(x\in\mathbb{R}\). What can you say about~\(f\)?
\def\ProblemHeadingText{Problem 3.}
\item
\begin{problemconcern}
\textbf{Warning: Suspected error.} The source statement was retained, but on a general finite measure space the displayed hypotheses do not imply the stated lower bound. An additional normalization such as \(\mu(X)=1\), or a different lower bound, may have been intended; the correction is not uniquely determined.
\end{problemconcern}
Let \(f\geq 0\) be~\(\mu\)-integrable such that \(\int f\,d\mu\leq 1\). Prove that \(\int \frac{1}{f}\,d\mu\geq 1\). (Hint, use the Schwarz inequality for \(\sqrt f\) and \(\frac{1}{\sqrt f}\).)
\def\ProblemHeadingText{Problem 4.}
\item
Let~\(R\) be a ring of subsets of~\(X\), \(E\) a Banach space and \(\mu:R\to E\) be a finitely additive measure. Prove that~\(\mu\) is~\(\sigma\)-additive iff for any decreasing sequence \((A_n)\) from~\(R\) with \(\bigcap_n A_n=\phi\) we have \(\lim \mu(A_n)=0\).
\end{examproblems}
\end{document}
